\documentclass[11pt,a4paper]{article}
\usepackage[utf8]{inputenc}
\usepackage[T1]{fontenc}
\usepackage[catalan]{babel}
\usepackage{lmodern}
\linespread{1.22}
\usepackage[a4paper,top=2.4cm,bottom=2.2cm,left=1.5cm,right=1.5cm,headheight=1.1cm,headsep=0.5cm,footskip=1cm]{geometry}
\usepackage{amsmath,amssymb,cancel}
\usepackage{icomma}
\usepackage{graphicx}
\usepackage[table]{xcolor}
\usepackage{tikz}\usetikzlibrary{calc}
\usepackage{pgfplots}\pgfplotsset{compat=1.17}
\usepackage{fontawesome5}
\usepackage{ccicons}
\usepackage{fancyhdr}
\usepackage{enumitem}
\usepackage{booktabs}
\usepackage{tabularx}
\usepackage{needspace}
\usepackage[most]{tcolorbox}
\usepackage{titlesec}
\usepackage[hidelinks]{hyperref}
\hypersetup{pdftitle={Apunts de teoria. Unitat 3: Sistemes d'equacions i inequacions. 1r de Batxillerat MI},pdfsubject={Matemàtiques I},pdfkeywords={sistemes d'equacions, mètode de Gauss, sistemes no lineals, sistemes d'inequacions, inequacions amb dues incògnites, semiplans}}

% colors de la web (FentMates, bloc d'Àlgebra: blau; es manté el nom fmgreen dels apunts de les altres unitats)
\definecolor{fmnavy}{HTML}{1B2748}
\definecolor{fmgreen}{HTML}{3D6AA0}
\definecolor{fmgreenl}{HTML}{4E81BD}
\definecolor{fmyellow}{HTML}{FFC94D}
\definecolor{fmgold}{HTML}{E9A919}
\definecolor{fmaccent}{HTML}{FF7A50}
\definecolor{fmteal}{HTML}{2BB3A3}
\definecolor{fmdef}{HTML}{1D4F73}\definecolor{fmdefbg}{HTML}{EEF4F9}
\definecolor{fmprop}{HTML}{3B4252}\definecolor{fmpropbg}{HTML}{F3F4F7}
\definecolor{fmex}{HTML}{A3203A}\definecolor{fmexbg}{HTML}{FBF0F2}
\definecolor{fmmuted}{HTML}{5D6685}\definecolor{fmball0}{HTML}{E2483D}\definecolor{fmball1}{HTML}{2F6FD6}
\definecolor{fmsurface}{HTML}{EEF1F8}
\definecolor{fmA}{HTML}{3B82F6}\definecolor{fmB}{HTML}{F97316}\definecolor{fmR}{HTML}{E11D48}\definecolor{fmK}{HTML}{16A34A}\definecolor{fmV}{HTML}{7C3AED}

\newcommand{\lt}{<}\newcommand{\gt}{>}
% fórmula en línia que no es parteix mai; si no cap en una línia, es posa centrada
\newcommand{\fmm}[1]{\setbox0=\hbox{$#1$}\ifdim\wd0>\linewidth\fmdisp{#1}\else\mbox{\box0}\fi}
% fórmula centrada; si és més ampla que la línia, s'encongeix fins que hi cap
\newcommand{\fmdisp}[1]{\[\setbox2=\hbox{$\displaystyle #1$}\ifdim\wd2>\linewidth\resizebox{\linewidth}{!}{\box2}\else\box2\fi\]}
\DeclareUnicodeCharacter{2212}{-}
\DeclareUnicodeCharacter{2225}{\ensuremath{\parallel}}\DeclareUnicodeCharacter{22A5}{\ensuremath{\perp}}\DeclareUnicodeCharacter{2226}{\ensuremath{\nparallel}}
\DeclareUnicodeCharacter{207A}{\ensuremath{^{+}}}\DeclareUnicodeCharacter{207B}{\ensuremath{^{-}}}
\setlength{\parindent}{0pt}
\setlength{\parskip}{12pt}
\setlength{\lineskip}{5pt}\setlength{\lineskiplimit}{5pt}  % més aire al voltant de les fraccions en línia
\setlength{\emergencystretch}{3em}
\setlist{itemsep=4pt,topsep=4pt}
\renewcommand{\thesection}{3.\arabic{section}}
\titleformat{\section}{\color{fmnavy}\Large\bfseries}{\colorbox{fmdef}{\color{white}\strut\ \thesection\ }}{0.6em}{}[{\color{fmdef}\titlerule[1.2pt]}]
\titlespacing*{\section}{0pt}{10pt}{4pt}
\titleformat*{\subsection}{\color{fmgreen}\large\bfseries}
\titlespacing*{\subsection}{0pt}{6pt plus 2pt}{0pt}
\allowdisplaybreaks

% dibuix de la targeta de la unitat a la web (illustrations.js), en blanc
\newcommand{\fmcard}{\begin{tikzpicture}[x=0.062cm,y=-0.062cm,white,line width=4pt,line cap=round]
  \begin{scope}[transparency group,opacity=0.15]
  \draw (20,106)--(220,20);
  \draw[opacity=0.75] (20,26)--(220,96);
  \draw[opacity=0.55] (128,6)--(128,114);
  \fill (128,59.5) circle (7);
  \node[anchor=base west,inner sep=0pt,font=\fontsize{21}{21}\selectfont] at (18,124) {$x+y+z=6$};
  \node[anchor=base west,inner sep=0pt,font=\fontsize{21}{21}\selectfont] at (150,12) {$2x-y=1$};
  \end{scope}
\end{tikzpicture}}
\newcommand{\fmbrand}{\raisebox{-0.25\height}{\includegraphics[height=0.5cm]{logo-fentmates.pdf}}}

% capçalera i peu
\pagestyle{fancy}\fancyhf{}
\fancyhead[L]{\small\fmbrand}
\fancyhead[R]{\small\color{fmmuted}1r Batxillerat MI · Unitat 3: Sistemes d'equacions i inequacions}
\fancyfoot[L]{\scriptsize\color{fmmuted}\href{https://creativecommons.org/licenses/by-nc-sa/4.0/deed.ca}{\ccbyncsa}\ \ FentMates · fentmates.org · CC BY-NC-SA 4.0}
\fancyfoot[R]{\small\color{fmmuted}\thepage}
\renewcommand{\headrulewidth}{0.4pt}\renewcommand{\headrule}{\color{fmmuted!40}\hrule height \headrulewidth}
\fancypagestyle{primera}{\fancyhf{}\renewcommand{\headrulewidth}{0pt}
  \fancyfoot[L]{\scriptsize\color{fmmuted}\href{https://creativecommons.org/licenses/by-nc-sa/4.0/deed.ca}{\ccbyncsa}\ \ FentMates · fentmates.org · CC BY-NC-SA 4.0}
  \fancyfoot[R]{\small\color{fmmuted}\thepage}}

% caixes (mateixos colors que la web)
% quan una caixa es parteix entre dues pàgines, la vora del tall és discontínua
\tcbset{fmbase/.style={enhanced,breakable,arc=3mm,boxrule=1pt,left=3.5mm,right=3.5mm,top=2mm,bottom=2.5mm,before skip=10pt,after skip=10pt,fonttitle=\bfseries,
  frame style={fill=tcbcolframe},
  overlay first={\fill[tcbcolback] ([xshift=1pt]frame.south west) rectangle ([xshift=-1pt,yshift=1.2pt]frame.south east);
    \draw[tcbcolframe,line width=1pt,dash pattern=on 5pt off 3pt] (frame.south west)--(frame.south east);},
  overlay middle={\fill[tcbcolback] ([xshift=1pt]frame.south west) rectangle ([xshift=-1pt,yshift=1.2pt]frame.south east);
    \fill[tcbcolback] ([xshift=1pt]frame.north west) rectangle ([xshift=-1pt,yshift=-1.2pt]frame.north east);
    \draw[tcbcolframe,line width=1pt,dash pattern=on 5pt off 3pt] (frame.south west)--(frame.south east) (frame.north west)--(frame.north east);},
  overlay last={\fill[tcbcolback] ([xshift=1pt]frame.north west) rectangle ([xshift=-1pt,yshift=-1.2pt]frame.north east);
    \draw[tcbcolframe,line width=1pt,dash pattern=on 5pt off 3pt] (frame.north west)--(frame.north east);}}}
\newtcolorbox{fmdef}[1]{fmbase,colback=fmdefbg,colframe=fmdef,coltitle=white,title={#1}}
\newtcolorbox{fmprop}[1]{fmbase,colback=fmpropbg,colframe=fmprop,coltitle=white,title={#1},halign title=center}
\newtcolorbox{fmex}[1]{fmbase,before upper=\raggedright,colback=fmexbg,colframe=fmex,colbacktitle=fmexbg,coltitle=fmex,title={\faPencil*\ #1},
  titlerule=0.4pt,titlerule style={fmex!40,dashed}}
\newtcolorbox{fmwarn}[1]{fmbase,colback=fmsf!6,colframe=fmsf,coltitle=white,title={#1}}
\newtcolorbox{fmnote}[1][]{enhanced,breakable,colback=white,colframe=white,borderline west={3pt}{0pt}{fmaccent},arc=0pt,left=3mm,right=2mm,top=1mm,bottom=1mm,before skip=6pt,after skip=6pt,before upper=\raggedright,#1}
\newcommand{\fmkey}[1]{\colorbox{fmyellow!45}{\strut #1}}
\newcommand{\fmuline}[1]{\textbf{\color{fmaccent}#1}}
\definecolor{fmok}{HTML}{15803D}\definecolor{fmev}{HTML}{2563EB}\definecolor{fmsf}{HTML}{C2410C}\definecolor{fmsi}{HTML}{7C3AED}\definecolor{fmas}{HTML}{A3203A}
\newcommand{\fmbadge}[2]{\tcbox[on line,colback=#1,colframe=#1,arc=2mm,boxsep=0pt,left=4pt,right=4pt,top=1.5pt,bottom=1.5pt]{\color{white}\footnotesize\bfseries #2}}
\newcommand{\fmchip}[1]{\tcbox[on line,colback=fmsf!10,colframe=fmsf!50,arc=2mm,boxsep=1pt,left=4pt,right=4pt]{#1}}
\newcommand{\fmlevel}[2]{\tcbox[on line,colback=white,colframe=fmmuted!40,arc=2mm,boxsep=2pt,halign=center]{\begin{tabular}{c}{\scriptsize\color{fmmuted}\MakeUppercase{#1}}\\#2\end{tabular}}}
% gràfiques
\pgfplotsset{fmaxis/.style={axis lines=middle,grid=major,grid style={fmmuted!15},axis line style={fmmuted,-latex},
  tick label style={font=\scriptsize,color=fmmuted},clip=true,unbounded coords=jump,restrict y to domain=-200:200,
  every axis plot/.append style={line join=round}},
  fmcurve/.style={fmaccent,line width=1.6pt,smooth}}
\tikzset{fmasy/.style={fmmuted,dashed,line width=0.9pt},fmguide/.style={fmaccent,dotted,line width=0.8pt},
  fmdot/.style={circle,fill=fmaccent,inner sep=0pt,minimum size=5pt},fmopen/.style={circle,draw=fmaccent,fill=white,line width=1pt,inner sep=0pt,minimum size=5pt}}
\newcommand{\fmpau}[1]{\tcbox[on line,colback=fmnavy,colframe=fmnavy,arc=2mm,boxsep=0pt,left=4pt,right=4pt,top=1.5pt,bottom=1.5pt]{\color{white}\scriptsize\bfseries\MakeUppercase{#1}}\ }

\begin{document}
\thispagestyle{primera}
\begin{tikzpicture}[remember picture,overlay]
  \fill[fmgreen] (current page.north west) rectangle ([yshift=-10.1cm]current page.north east);
  \begin{scope}
    \clip (current page.north west) rectangle ([yshift=-10.1cm]current page.north east);
    \node[anchor=north west,inner sep=0pt] at ([xshift=6.6cm,yshift=-1.3cm]current page.north west) {\fmcard};
  \end{scope}
  \node[anchor=north west,inner sep=0pt] at ([xshift=1.5cm,yshift=-1.1cm]current page.north west) {\tcbox[colback=white,colframe=white,arc=2.5mm,boxsep=0pt,left=7pt,right=7pt,top=5pt,bottom=5pt]{\includegraphics[height=0.75cm]{logo-fentmates.pdf}}};
\end{tikzpicture}
\vspace*{0.0cm}
\tcbox[colback=white!15!fmgreen,colframe=white!15!fmgreen,arc=2.5mm,boxsep=0pt,left=6pt,right=6pt,top=3pt,bottom=3pt]{\color{white}\small\bfseries Apunts de teoria · 1r Batxillerat · Matemàtiques I}

\vspace{4pt}
{\raggedright\color{white}\fontsize{25}{30}\selectfont\bfseries{\color{fmyellow}Unitat 3} · Sistemes d'equacions i inequacions\par}
\vspace{6pt}
\begin{minipage}{\linewidth}\color{white!88!fmgreen}\small Teoria i exemples resolts pas a pas: els sistemes lineals pel mètode de Gauss i els no lineals, els sistemes d'inequacions amb una incògnita, les inequacions i els sistemes amb dues incògnites i la resolució de problemes.\end{minipage}\par
% l'índex comença sempre a la mateixa altura, per davall de la franja (com en les unitats 5-10)
\vspace*{\dimexpr 8.6cm-\pagetotal\relax}

\begin{tcolorbox}[colback=fmsurface,colframe=fmsurface,arc=3mm,left=4mm,right=4mm,top=2mm,bottom=2mm]
\small\textbf{Índex}\par\smallskip
\begin{minipage}[t]{0.49\linewidth}\small \mbox{\hyperlink{sec1}{\textbf{\color{fmgreen}3.1}\ \ Sistemes lineals. Gauss}}\\[3pt]
\mbox{\hyperlink{sec2}{\textbf{\color{fmgreen}3.2}\ \ Sistemes no lineals}}\\[3pt]
\hyperlink{sec3}{\textbf{\color{fmgreen}3.3}\ \ Sistemes d'inequacions amb una incògnita}\\[3pt]\end{minipage}\hfill\begin{minipage}[t]{0.49\linewidth}\small \mbox{\hyperlink{sec4}{\textbf{\color{fmgreen}3.4}\ \ Inequacions amb dues incògnites}}\\[3pt]
\hyperlink{sec5}{\textbf{\color{fmgreen}3.5}\ \ Sistemes d'inequacions amb dues incògnites}\\[3pt]
\mbox{\hyperlink{sec6}{\textbf{\color{fmgreen}3.6}\ \ Resolució de problemes}}\\[3pt]\end{minipage}
\end{tcolorbox}

\Needspace{8\baselineskip}
\section{Sistemes d'equacions lineals. Mètode de Gauss}\hypertarget{sec1}{}
\begin{fmdef}{Definicions}
\begin{itemize}
\item Un \textbf{sistema d'equacions lineals} és un conjunt d'equacions de primer grau amb diverses incògnites. Una \textbf{solució} és un conjunt de valors de les incògnites que verifica totes les equacions alhora.
\item Dos sistemes són \textbf{equivalents} si tenen les mateixes solucions. N'obtenim un d'equivalent si intercanviem dues equacions, multipliquem una equació per un nombre \fmm{\neq0} o sumem a una equació un múltiple d'una altra.
\end{itemize}

\begin{center}\renewcommand{\arraystretch}{1.5}\setlength{\parskip}{0pt}\small\begin{tabularx}{\linewidth}{>{\raggedright\arraybackslash}p{3.2cm}>{\raggedright\arraybackslash}X>{\raggedright\arraybackslash}X}\toprule
\textbf{Tipus} & \textbf{Solucions} & \textbf{Significat}\\\midrule
\fmbadge{fmok}{SCD} Compatible determinat & una única solució & les rectes o plans es tallen en un punt\\[3pt]\end{tabularx}\par\nointerlineskip\vskip0pt
\small\begin{tabularx}{\linewidth}{>{\raggedright\arraybackslash}p{3.2cm}>{\raggedright\arraybackslash}X>{\raggedright\arraybackslash}X}\fmbadge{fmev}{SCI} Compatible indeterminat & infinites solucions & depenen d'un o més paràmetres\\[3pt]\end{tabularx}\par\nointerlineskip\vskip0pt
\small\begin{tabularx}{\linewidth}{>{\raggedright\arraybackslash}p{3.2cm}>{\raggedright\arraybackslash}X>{\raggedright\arraybackslash}X}\fmbadge{fmas}{SI} Incompatible & cap solució & equacions contradictòries\\[3pt]\bottomrule\end{tabularx}\end{center}
\end{fmdef}

\Needspace{5\baselineskip}
\begin{fmprop}{Mètode de Gauss}
Consisteix a transformar el sistema en un altre d'equivalent \textbf{esglaonat} (triangular), en què cada equació té una incògnita menys que l'anterior. Treballem amb la \textbf{matriu ampliada} (coeficients i termes independents):

\fmdisp{\left(\begin{array}{ccc|c}a&b&c&d\\ a'&b'&c'&d'\\ a''&b''&c''&d''\end{array}\right)\longrightarrow\left(\begin{array}{ccc|c}*&*&*&*\\ 0&*&*&*\\ 0&0&m&n\end{array}\right)}

\begin{enumerate}
\item Fem zeros en la primera columna, davall del primer element, combinant cada fila amb la primera: \fmm{F_2\to aF_2-a'F_1}\dots{}
\item Fem zeros en la segona columna davall del segon element, combinant amb la segona fila.
\item Mirem l'última fila \fmm{\left(0\ \ 0\ \ m\,|\,n\right)}:

\begin{itemize}
\item \fmm{m\neq0}: \textbf{SCD}. S'aïlla \fmm{z=\frac nm} i se substitueix cap amunt.
\item \fmm{m=0,\ n=0}: \textbf{SCI}. L'última equació desapareix (\fmm{0=0}) i una incògnita passa a ser un paràmetre (\fmm{z=\lambda}).
\item \fmm{m=0,\ n\neq0}: \textbf{SI}. L'última equació és \fmm{0=n}, impossible.
\end{itemize}
\end{enumerate}

\begin{fmnote}\fmuline{Consell}: convé posar en primer lloc una equació amb coeficient 1 (o \fmm{-1}) en la primera incògnita.\end{fmnote}
\end{fmprop}

\Needspace{5\baselineskip}
\begin{fmex}{Exemple resolt: mètode de Gauss}
\begin{enumerate}[label=\textbf{\arabic*.},leftmargin=*]
\item Resol \fmm{\begin{cases}x+y+z=6\\ 2x-y+z=3\\ x+2y-z=2\end{cases}}

\fmm{\left(\begin{array}{rrr|r}1&1&1&6\\ 2&-1&1&3\\ 1&2&-1&2\end{array}\right)\xrightarrow[F_3\to F_3-F_1]{F_2\to F_2-2F_1}\left(\begin{array}{rrr|r}1&1&1&6\\ 0&-3&-1&-9\\ 0&1&-2&-4\end{array}\right)}

\fmm{\xrightarrow{F_3\to 3F_3+F_2}\left(\begin{array}{rrr|r}1&1&1&6\\ 0&-3&-1&-9\\ 0&0&-7&-21\end{array}\right)}

\fmm{-7z=-21\Rightarrow z=3;\quad -3y-3=-9\Rightarrow y=2;\quad x+2+3=6\Rightarrow x=1}. \textbf{SCD}: \fmm{(x,y,z)=(1,2,3)}
\item Resol \fmm{\begin{cases}x-y+z=1\\ 2x+y-z=5\\ x+2y-2z=4\end{cases}}

\fmm{\left(\begin{array}{rrr|r}1&-1&1&1\\ 2&1&-1&5\\ 1&2&-2&4\end{array}\right)\xrightarrow[F_3\to F_3-F_1]{F_2\to F_2-2F_1}\left(\begin{array}{rrr|r}1&-1&1&1\\ 0&3&-3&3\\ 0&3&-3&3\end{array}\right)\xrightarrow{F_3\to F_3-F_2}\left(\begin{array}{rrr|r}1&-1&1&1\\ 0&3&-3&3\\ 0&0&0&0\end{array}\right)}

L'última fila és \fmm{0=0}: \textbf{SCI}. Fem \fmm{z=\lambda}: \fmm{\ 3y-3\lambda=3\Rightarrow y=1+\lambda;\ \ x=1+y-z=2}

\fmm{(x,y,z)=(2,\ 1+\lambda,\ \lambda),\ \lambda\in\mathbb R}. Si el terme independent de la tercera equació fora 5, l'última fila quedaria \fmm{\left(0\ \ 0\ \ 0\,|\,1\right)}: \fmm{0=1}, el sistema seria \textbf{incompatible}.
\end{enumerate}
\end{fmex}

\Needspace{8\baselineskip}
\section{Sistemes d'equacions no lineals}\hypertarget{sec2}{}
\begin{fmdef}{Sistemes no lineals}
Són sistemes en què alguna equació no és de primer grau (apareixen productes d'incògnites, quadrats, arrels, logaritmes, exponencials\dots{}). Els mètodes més habituals són:

\begin{itemize}
\item \textbf{Substitució}: s'aïlla una incògnita (preferiblement d'una equació lineal) i se substitueix en l'altra.
\item \textbf{Reducció}: si les incògnites apareixen amb la mateixa forma (\fmm{x^2} i \fmm{y^2}, per exemple), es combinen les equacions per a eliminar-ne una.
\item \textbf{Canvi de variable}: en sistemes exponencials (\fmm{u=2^x,\ v=3^y}) o logarítmics (\fmm{u=\log x,\ v=\log y}) s'obté un sistema lineal en \fmm{u} i \fmm{v}.
\end{itemize}

\begin{fmnote}Sempre cal \textbf{comprovar} les solucions (denominadors, arrels, logaritmes\dots{}).\end{fmnote}
\end{fmdef}

\Needspace{5\baselineskip}
\begin{fmex}{Exemple resolt: sistemes no lineals}
\begin{enumerate}[label=\textbf{\arabic*.},leftmargin=*]
\item \fmm{\begin{cases}x-y=1\\ x^2+y^2=13\end{cases}}

\fmm{x=1+y\Rightarrow(1+y)^2+y^2=13\Rightarrow 2y^2+2y-12=0\Rightarrow y^2+y-6=0}

\fmm{y=2\Rightarrow x=3;\quad y=-3\Rightarrow x=-2}. Solucions: \fmm{(3,2)} i \fmm{(-2,-3)}
\item \fmm{\begin{cases}2^x+3^y=7\\ 2^{x+1}-3^y=5\end{cases}}

\fmm{u=2^x,\ v=3^y:\ \ u+v=7,\ \ 2u-v=5}. Sumant: \fmm{3u=12\Rightarrow u=4,\ v=3}

\fmm{2^x=4\Rightarrow x=2;\quad 3^y=3\Rightarrow y=1}
\item \fmm{\begin{cases}\log x+\log y=3\\ x-y=90\end{cases}}

\fmm{\log(xy)=\log1000\Rightarrow xy=1000}. Amb \fmm{x=90+y:\ \ y^2+90y-1000=0}

\fmm{y=10} o \fmm{y=-100} (no vàlida). Solució: \fmm{x=100,\ y=10}
\end{enumerate}
\end{fmex}

\Needspace{8\baselineskip}
\section{Sistemes d'inequacions amb una incògnita}\hypertarget{sec3}{}
\begin{fmprop}{Com resoldre'ls}
\begin{itemize}
\item Es resol cada inequació per separat, com a la unitat 2, i la solució del sistema és la \textbf{intersecció} de totes.
\item Convé representar la solució de cada inequació en una recta real, una sobre l'altra i amb la mateixa escala, per a veure què tenen en comú. Si no tenen cap punt en comú, el sistema \textbf{no té solució}.
\end{itemize}
\end{fmprop}

\Needspace{5\baselineskip}
\begin{fmex}{Exemple resolt: sistemes d'inequacions}
\begin{enumerate}[label=\textbf{\arabic*.},leftmargin=*]
\item Resol el sistema \fmm{\begin{cases}2x-1\lt5\\ x^2-4\ge0\end{cases}}

\fmm{2x\lt6\Rightarrow x\lt3:\ (-\infty,3).\qquad x^2-4\ge0\Rightarrow(-\infty,-2]\cup[2,+\infty)}

Representem la solució de cada inequació en una recta real, una sobre l'altra:

\begin{center}\begin{tikzpicture}[line cap=round]\draw[fmmuted!60,dashed,line width=0.5pt] (4.678,0.1)--(4.678,-1.700);\draw[fmmuted!60,dashed,line width=0.5pt] (9.433,0.1)--(9.433,-1.700);\draw[fmmuted!60,dashed,line width=0.5pt] (10.622,0.1)--(10.622,-1.700);\node[anchor=west,font=\small] at (0,-0.400) {$2x-1<5$};\draw[fmmuted!50,line width=0.6pt] (2.300,-0.400)--(13.000,-0.400);\draw[fmA,line width=2.6pt] (2.300,-0.400)--(10.622,-0.400);\fill[fmA] (2.050,-0.400)--++(0.3,0.19)--++(0,-0.38)--cycle;\draw[fmA,line width=1.2pt,fill=white] (10.622,-0.400) circle[radius=0.13];\node[anchor=west,font=\small] at (0,-1.150) {$x^2-4\geq0$};\draw[fmmuted!50,line width=0.6pt] (2.300,-1.150)--(13.000,-1.150);\draw[fmB,line width=2.6pt] (2.300,-1.150)--(4.678,-1.150);\fill[fmB] (2.050,-1.150)--++(0.3,0.19)--++(0,-0.38)--cycle;\draw[fmB,line width=1.2pt,fill=fmB] (4.678,-1.150) circle[radius=0.13];\draw[fmB,line width=2.6pt] (9.433,-1.150)--(13.000,-1.150);\fill[fmB] (13.250,-1.150)--++(-0.3,0.19)--++(0,-0.38)--cycle;\draw[fmB,line width=1.2pt,fill=fmB] (9.433,-1.150) circle[radius=0.13];\draw[fmmuted,line width=0.8pt,-latex] (2.100,-1.700)--(13.400,-1.700);\draw[fmmuted,line width=0.6pt] (2.300,-1.700)--++(0,0.12)--++(0,-0.24) node[below,font=\scriptsize,text=fmmuted] {$-4$};\draw[fmmuted,line width=0.6pt] (3.489,-1.700)--++(0,0.12)--++(0,-0.24) node[below,font=\scriptsize,text=fmmuted] {$-3$};\draw[fmmuted,line width=0.6pt] (4.678,-1.700)--++(0,0.12)--++(0,-0.24) node[below,font=\scriptsize,text=fmmuted] {$-2$};\draw[fmmuted,line width=0.6pt] (5.867,-1.700)--++(0,0.12)--++(0,-0.24) node[below,font=\scriptsize,text=fmmuted] {$-1$};\draw[fmmuted,line width=0.6pt] (7.056,-1.700)--++(0,0.12)--++(0,-0.24) node[below,font=\scriptsize,text=fmmuted] {$0$};\draw[fmmuted,line width=0.6pt] (8.244,-1.700)--++(0,0.12)--++(0,-0.24) node[below,font=\scriptsize,text=fmmuted] {$1$};\draw[fmmuted,line width=0.6pt] (9.433,-1.700)--++(0,0.12)--++(0,-0.24) node[below,font=\scriptsize,text=fmmuted] {$2$};\draw[fmmuted,line width=0.6pt] (10.622,-1.700)--++(0,0.12)--++(0,-0.24) node[below,font=\scriptsize,text=fmmuted] {$3$};\draw[fmmuted,line width=0.6pt] (11.811,-1.700)--++(0,0.12)--++(0,-0.24) node[below,font=\scriptsize,text=fmmuted] {$4$};\draw[fmmuted,line width=0.6pt] (13.000,-1.700)--++(0,0.12)--++(0,-0.24) node[below,font=\scriptsize,text=fmmuted] {$5$};\end{tikzpicture}\end{center}

La solució del sistema és la part on coincideixen les dues rectes, la \textbf{intersecció}: \fmm{x\in(-\infty,-2]\cup[2,3)}

\begin{center}\begin{tikzpicture}[line cap=round]\fill[fmR,opacity=0.12] (2.300,0.1) rectangle (4.678,-2.450);\fill[fmR,opacity=0.12] (9.433,0.1) rectangle (10.622,-2.450);\draw[fmmuted!60,dashed,line width=0.5pt] (4.678,0.1)--(4.678,-2.450);\draw[fmmuted!60,dashed,line width=0.5pt] (9.433,0.1)--(9.433,-2.450);\draw[fmmuted!60,dashed,line width=0.5pt] (10.622,0.1)--(10.622,-2.450);\node[anchor=west,font=\small] at (0,-0.400) {$2x-1<5$};\draw[fmmuted!50,line width=0.6pt] (2.300,-0.400)--(13.000,-0.400);\draw[fmA,line width=2.6pt] (2.300,-0.400)--(10.622,-0.400);\fill[fmA] (2.050,-0.400)--++(0.3,0.19)--++(0,-0.38)--cycle;\draw[fmA,line width=1.2pt,fill=white] (10.622,-0.400) circle[radius=0.13];\node[anchor=west,font=\small] at (0,-1.150) {$x^2-4\geq0$};\draw[fmmuted!50,line width=0.6pt] (2.300,-1.150)--(13.000,-1.150);\draw[fmB,line width=2.6pt] (2.300,-1.150)--(4.678,-1.150);\fill[fmB] (2.050,-1.150)--++(0.3,0.19)--++(0,-0.38)--cycle;\draw[fmB,line width=1.2pt,fill=fmB] (4.678,-1.150) circle[radius=0.13];\draw[fmB,line width=2.6pt] (9.433,-1.150)--(13.000,-1.150);\fill[fmB] (13.250,-1.150)--++(-0.3,0.19)--++(0,-0.38)--cycle;\draw[fmB,line width=1.2pt,fill=fmB] (9.433,-1.150) circle[radius=0.13];\node[anchor=west,font=\small] at (0,-1.900) {$\text{Solució}$};\draw[fmmuted!50,line width=0.6pt] (2.300,-1.900)--(13.000,-1.900);\draw[fmR,line width=2.6pt] (2.300,-1.900)--(4.678,-1.900);\fill[fmR] (2.050,-1.900)--++(0.3,0.19)--++(0,-0.38)--cycle;\draw[fmR,line width=1.2pt,fill=fmR] (4.678,-1.900) circle[radius=0.13];\draw[fmR,line width=2.6pt] (9.433,-1.900)--(10.622,-1.900);\draw[fmR,line width=1.2pt,fill=fmR] (9.433,-1.900) circle[radius=0.13];\draw[fmR,line width=1.2pt,fill=white] (10.622,-1.900) circle[radius=0.13];\draw[fmmuted,line width=0.8pt,-latex] (2.100,-2.450)--(13.400,-2.450);\draw[fmmuted,line width=0.6pt] (2.300,-2.450)--++(0,0.12)--++(0,-0.24) node[below,font=\scriptsize,text=fmmuted] {$-4$};\draw[fmmuted,line width=0.6pt] (3.489,-2.450)--++(0,0.12)--++(0,-0.24) node[below,font=\scriptsize,text=fmmuted] {$-3$};\draw[fmmuted,line width=0.6pt] (4.678,-2.450)--++(0,0.12)--++(0,-0.24) node[below,font=\scriptsize,text=fmmuted] {$-2$};\draw[fmmuted,line width=0.6pt] (5.867,-2.450)--++(0,0.12)--++(0,-0.24) node[below,font=\scriptsize,text=fmmuted] {$-1$};\draw[fmmuted,line width=0.6pt] (7.056,-2.450)--++(0,0.12)--++(0,-0.24) node[below,font=\scriptsize,text=fmmuted] {$0$};\draw[fmmuted,line width=0.6pt] (8.244,-2.450)--++(0,0.12)--++(0,-0.24) node[below,font=\scriptsize,text=fmmuted] {$1$};\draw[fmmuted,line width=0.6pt] (9.433,-2.450)--++(0,0.12)--++(0,-0.24) node[below,font=\scriptsize,text=fmmuted] {$2$};\draw[fmmuted,line width=0.6pt] (10.622,-2.450)--++(0,0.12)--++(0,-0.24) node[below,font=\scriptsize,text=fmmuted] {$3$};\draw[fmmuted,line width=0.6pt] (11.811,-2.450)--++(0,0.12)--++(0,-0.24) node[below,font=\scriptsize,text=fmmuted] {$4$};\draw[fmmuted,line width=0.6pt] (13.000,-2.450)--++(0,0.12)--++(0,-0.24) node[below,font=\scriptsize,text=fmmuted] {$5$};\end{tikzpicture}\end{center}

\begin{fmnote}El 3 no s'inclou perquè no és solució de la primera inequació. El \fmm{-2} i el 2 sí, perquè són solució de totes dues.\end{fmnote}
\item Resol el sistema \fmm{\begin{cases}x-2\gt0\\ x+1\lt0\end{cases}}

\fmm{x\gt2:\ (2,+\infty).\qquad x\lt-1:\ (-\infty,-1)}

Representem la solució de cada inequació en una recta real, una sobre l'altra:

\begin{center}\begin{tikzpicture}[line cap=round]\draw[fmmuted!60,dashed,line width=0.5pt] (5.357,0.1)--(5.357,-1.700);\draw[fmmuted!60,dashed,line width=0.5pt] (9.943,0.1)--(9.943,-1.700);\node[anchor=west,font=\small] at (0,-0.400) {$x-2>0$};\draw[fmmuted!50,line width=0.6pt] (2.300,-0.400)--(13.000,-0.400);\draw[fmA,line width=2.6pt] (9.943,-0.400)--(13.000,-0.400);\fill[fmA] (13.250,-0.400)--++(-0.3,0.19)--++(0,-0.38)--cycle;\draw[fmA,line width=1.2pt,fill=white] (9.943,-0.400) circle[radius=0.13];\node[anchor=west,font=\small] at (0,-1.150) {$x+1<0$};\draw[fmmuted!50,line width=0.6pt] (2.300,-1.150)--(13.000,-1.150);\draw[fmB,line width=2.6pt] (2.300,-1.150)--(5.357,-1.150);\fill[fmB] (2.050,-1.150)--++(0.3,0.19)--++(0,-0.38)--cycle;\draw[fmB,line width=1.2pt,fill=white] (5.357,-1.150) circle[radius=0.13];\draw[fmmuted,line width=0.8pt,-latex] (2.100,-1.700)--(13.400,-1.700);\draw[fmmuted,line width=0.6pt] (2.300,-1.700)--++(0,0.12)--++(0,-0.24) node[below,font=\scriptsize,text=fmmuted] {$-3$};\draw[fmmuted,line width=0.6pt] (3.829,-1.700)--++(0,0.12)--++(0,-0.24) node[below,font=\scriptsize,text=fmmuted] {$-2$};\draw[fmmuted,line width=0.6pt] (5.357,-1.700)--++(0,0.12)--++(0,-0.24) node[below,font=\scriptsize,text=fmmuted] {$-1$};\draw[fmmuted,line width=0.6pt] (6.886,-1.700)--++(0,0.12)--++(0,-0.24) node[below,font=\scriptsize,text=fmmuted] {$0$};\draw[fmmuted,line width=0.6pt] (8.414,-1.700)--++(0,0.12)--++(0,-0.24) node[below,font=\scriptsize,text=fmmuted] {$1$};\draw[fmmuted,line width=0.6pt] (9.943,-1.700)--++(0,0.12)--++(0,-0.24) node[below,font=\scriptsize,text=fmmuted] {$2$};\draw[fmmuted,line width=0.6pt] (11.471,-1.700)--++(0,0.12)--++(0,-0.24) node[below,font=\scriptsize,text=fmmuted] {$3$};\draw[fmmuted,line width=0.6pt] (13.000,-1.700)--++(0,0.12)--++(0,-0.24) node[below,font=\scriptsize,text=fmmuted] {$4$};\end{tikzpicture}\end{center}

Els dos intervals no tenen cap punt en comú: la intersecció és buida i el sistema no té solució.

\begin{center}\begin{tikzpicture}[line cap=round]\draw[fmmuted!60,dashed,line width=0.5pt] (5.357,0.1)--(5.357,-2.450);\draw[fmmuted!60,dashed,line width=0.5pt] (9.943,0.1)--(9.943,-2.450);\node[anchor=west,font=\small] at (0,-0.400) {$x-2>0$};\draw[fmmuted!50,line width=0.6pt] (2.300,-0.400)--(13.000,-0.400);\draw[fmA,line width=2.6pt] (9.943,-0.400)--(13.000,-0.400);\fill[fmA] (13.250,-0.400)--++(-0.3,0.19)--++(0,-0.38)--cycle;\draw[fmA,line width=1.2pt,fill=white] (9.943,-0.400) circle[radius=0.13];\node[anchor=west,font=\small] at (0,-1.150) {$x+1<0$};\draw[fmmuted!50,line width=0.6pt] (2.300,-1.150)--(13.000,-1.150);\draw[fmB,line width=2.6pt] (2.300,-1.150)--(5.357,-1.150);\fill[fmB] (2.050,-1.150)--++(0.3,0.19)--++(0,-0.38)--cycle;\draw[fmB,line width=1.2pt,fill=white] (5.357,-1.150) circle[radius=0.13];\node[anchor=west,font=\small] at (0,-1.900) {$\text{Solució}$};\draw[fmmuted!50,line width=0.6pt] (2.300,-1.900)--(13.000,-1.900);\node[font=\small] at (7.650,-1.650) {$\emptyset$};\draw[fmmuted,line width=0.8pt,-latex] (2.100,-2.450)--(13.400,-2.450);\draw[fmmuted,line width=0.6pt] (2.300,-2.450)--++(0,0.12)--++(0,-0.24) node[below,font=\scriptsize,text=fmmuted] {$-3$};\draw[fmmuted,line width=0.6pt] (3.829,-2.450)--++(0,0.12)--++(0,-0.24) node[below,font=\scriptsize,text=fmmuted] {$-2$};\draw[fmmuted,line width=0.6pt] (5.357,-2.450)--++(0,0.12)--++(0,-0.24) node[below,font=\scriptsize,text=fmmuted] {$-1$};\draw[fmmuted,line width=0.6pt] (6.886,-2.450)--++(0,0.12)--++(0,-0.24) node[below,font=\scriptsize,text=fmmuted] {$0$};\draw[fmmuted,line width=0.6pt] (8.414,-2.450)--++(0,0.12)--++(0,-0.24) node[below,font=\scriptsize,text=fmmuted] {$1$};\draw[fmmuted,line width=0.6pt] (9.943,-2.450)--++(0,0.12)--++(0,-0.24) node[below,font=\scriptsize,text=fmmuted] {$2$};\draw[fmmuted,line width=0.6pt] (11.471,-2.450)--++(0,0.12)--++(0,-0.24) node[below,font=\scriptsize,text=fmmuted] {$3$};\draw[fmmuted,line width=0.6pt] (13.000,-2.450)--++(0,0.12)--++(0,-0.24) node[below,font=\scriptsize,text=fmmuted] {$4$};\end{tikzpicture}\end{center}
\end{enumerate}
\end{fmex}

\Needspace{8\baselineskip}
\section{Inequacions amb dues incògnites}\hypertarget{sec4}{}
\begin{fmdef}{Inequacions lineals amb dues incògnites}
Una inequació lineal amb dues incògnites té la forma \fmm{ax+by\lt c} (o amb \fmm{\gt,\le,\ge}). Una \textbf{solució} és un parell de nombres \fmm{(x,y)} que la compleix, és a dir, un punt del pla. N'hi ha infinites: totes juntes formen un \textbf{semiplà}, una de les dues regions en què la recta \fmm{ax+by=c} divideix el pla.
\end{fmdef}

\Needspace{5\baselineskip}
\begin{fmprop}{Com es resolen, pas a pas}
\begin{enumerate}
\item \textbf{Recta frontera.} Canviem la desigualtat per un igual: \fmm{ax+by=c}.
\item \textbf{Dibuixem la recta} a partir de dos punts (per exemple, els talls amb els eixos, fent \fmm{x=0} i \fmm{y=0}). És \textbf{contínua} si la desigualtat és \fmm{\le} o \fmm{\ge}, perquè els punts de la recta són solució, i \textbf{discontínua} si és \fmm{\lt} o \fmm{\gt}, perquè no ho són.
\item \textbf{Punt de prova.} Triem un punt que no estiga en la recta (sovint \fmm{O(0,0)}; si la recta passa per l'origen, un altre, com \fmm{(1,0)}) i el substituïm en la inequació.
\item \textbf{Triem el semiplà.} Si el punt compleix la inequació, la solució és el semiplà que el conté; si no, és l'altre. L'ombregem.
\end{enumerate}
\end{fmprop}

\Needspace{5\baselineskip}
\begin{fmex}{Exemple resolt: inequacions amb dues incògnites}
\begin{enumerate}[label=\textbf{\arabic*.},leftmargin=*]
\item Resol \fmm{x+2y\le6}.

Recta frontera: \fmm{x+2y=6}. Dos punts: \fmm{(6,0)} i \fmm{(0,3)}. Com que hi ha \fmm{\le}, la recta és contínua.

Punt de prova \fmm{O(0,0)}: \fmm{0+2\cdot0=0\le6\ \checkmark}

El punt compleix la inequació: la solució és el semiplà que conté l'origen, amb la recta inclosa.

\begin{center}\begin{tikzpicture}\begin{axis}[fmaxis,axis equal image,width=5.6cm,height=5.6cm,xmin=-3,xmax=7,ymin=-3,ymax=7,xtick={-2,2,4,6},ytick={-2,2,4,6},minor xtick={-3,-2,-1,0,1,2,3,4,5,6,7},minor ytick={-3,-2,-1,0,1,2,3,4,5,6,7},grid=both,minor grid style={fmmuted!15},major grid style={fmmuted!15},tick label style={font=\tiny}]\fill[fmR,opacity=0.22] (axis cs:-3.0000,-3.0000)--(axis cs:7.0000,-3.0000)--(axis cs:7.0000,-0.5000)--(axis cs:-3.0000,4.5000)--cycle;\draw[fmA,line width=1.5pt] (axis cs:-4.0000,5.0000)--(axis cs:8.0000,-1.0000);\node[circle,fill=fmK,inner sep=0pt,minimum size=6pt] at (axis cs:0,0) {};\end{axis}\end{tikzpicture}\end{center}
\item Resol \fmm{y\gt x+1}.

Recta frontera: \fmm{y=x+1}. Dos punts: \fmm{(-1,0)} i \fmm{(0,1)}. Com que hi ha \fmm{\gt}, la recta és discontínua.

Punt de prova \fmm{O(0,0)}: \fmm{0\gt0+1} és fals.

El punt no compleix la inequació: la solució és el semiplà que no conté l'origen, sense la recta.

\begin{center}\begin{tikzpicture}\begin{axis}[fmaxis,axis equal image,width=5.6cm,height=5.6cm,xmin=-5,xmax=5,ymin=-5,ymax=5,xtick={-4,-2,2,4},ytick={-4,-2,2,4},minor xtick={-5,-4,-3,-2,-1,0,1,2,3,4,5},minor ytick={-5,-4,-3,-2,-1,0,1,2,3,4,5},grid=both,minor grid style={fmmuted!15},major grid style={fmmuted!15},tick label style={font=\tiny}]\fill[fmR,opacity=0.22] (axis cs:4.0000,5.0000)--(axis cs:-5.0000,5.0000)--(axis cs:-5.0000,-4.0000)--cycle;\draw[fmA,line width=1.5pt,dash pattern=on 5pt off 3.5pt] (axis cs:-6.0000,-5.0000)--(axis cs:6.0000,7.0000);\node[circle,fill=fmR,inner sep=0pt,minimum size=6pt] at (axis cs:0,0) {};\end{axis}\end{tikzpicture}\end{center}
\item Resol \fmm{y\ge2x}.

Recta frontera: \fmm{y=2x}. Dos punts: \fmm{(0,0)} i \fmm{(1,2)}. Recta contínua.

La recta passa per l'origen, així que hem de provar un altre punt: \fmm{(1,0)}: \fmm{0\ge2\cdot1} és fals.

La solució és el semiplà que no conté el punt \fmm{(1,0)}, amb la recta inclosa.

\begin{center}\begin{tikzpicture}\begin{axis}[fmaxis,axis equal image,width=5.6cm,height=5.6cm,xmin=-5,xmax=5,ymin=-5,ymax=5,xtick={-4,-2,2,4},ytick={-4,-2,2,4},minor xtick={-5,-4,-3,-2,-1,0,1,2,3,4,5},minor ytick={-5,-4,-3,-2,-1,0,1,2,3,4,5},grid=both,minor grid style={fmmuted!15},major grid style={fmmuted!15},tick label style={font=\tiny}]\fill[fmR,opacity=0.22] (axis cs:-5.0000,-5.0000)--(axis cs:-2.5000,-5.0000)--(axis cs:2.5000,5.0000)--(axis cs:-5.0000,5.0000)--cycle;\draw[fmA,line width=1.5pt] (axis cs:-6.0000,-12.0000)--(axis cs:6.0000,12.0000);\node[circle,fill=fmR,inner sep=0pt,minimum size=6pt] at (axis cs:1,0) {};\end{axis}\end{tikzpicture}\end{center}
\end{enumerate}
\end{fmex}

\Needspace{8\baselineskip}
\section{Sistemes d'inequacions amb dues incògnites}\hypertarget{sec5}{}
\begin{fmprop}{Com es resolen}
\begin{enumerate}
\item Es resol cada inequació per separat, com a l'apartat anterior: cada una dona un semiplà.
\item La solució del sistema és la \textbf{intersecció} de tots els semiplans, la \textbf{regió factible}. Si no tenen cap punt en comú, el sistema no té solució.
\item Els \textbf{vèrtexs} de la regió són els punts de tall de les rectes frontera: es calculen resolent el sistema format per les dues equacions.
\item La regió pot ser \textbf{fitada} (un polígon) o \textbf{no fitada} (s'estén indefinidament).
\end{enumerate}
\end{fmprop}

\Needspace{5\baselineskip}
\begin{fmex}{Exemple resolt: sistemes amb dues incògnites}
\begin{enumerate}[label=\textbf{\arabic*.},leftmargin=*]
\item Resol \fmm{\begin{cases}x+y\le4\\ y\ge x-2\\ x\ge0,\ y\ge0\end{cases}}

\fmm{x+y\le4}: recta \fmm{x+y=4}; amb \fmm{O(0,0)}: \fmm{0\le4\ \checkmark}: semiplà de \fmm{O}.

\fmm{y\ge x-2}: recta \fmm{y=x-2}; amb \fmm{O}: \fmm{0\ge-2\ \checkmark}: semiplà de \fmm{O}. \fmm{\ x\ge0,\ y\ge0}: primer quadrant.

La solució és el polígon de vèrtexs \fmm{(0,0),\ (2,0),\ (3,1)} i \fmm{(0,4)}. La regió és fitada.

\begin{center}\begin{tikzpicture}\begin{axis}[fmaxis,axis equal image,width=5.6cm,height=5.6cm,xmin=-2,xmax=6,ymin=-2,ymax=6,xtick={-2,2,4,6},ytick={-2,2,4,6},minor xtick={-2,-1,0,1,2,3,4,5,6},minor ytick={-2,-1,0,1,2,3,4,5,6},grid=both,minor grid style={fmmuted!15},major grid style={fmmuted!15},tick label style={font=\tiny}]\fill[fmR,opacity=0.22] (axis cs:2.0000,0.0000)--(axis cs:3.0000,1.0000)--(axis cs:0.0000,4.0000)--(axis cs:0.0000,0.0000)--cycle;\draw[fmA,line width=1.5pt] (axis cs:-3.0000,7.0000)--(axis cs:7.0000,-3.0000);\draw[fmB,line width=1.5pt] (axis cs:-3.0000,-5.0000)--(axis cs:7.0000,5.0000);\draw[fmV,line width=1.5pt] (axis cs:0.0000,-3.0000)--(axis cs:0.0000,7.0000);\draw[fmK,line width=1.5pt] (axis cs:-3.0000,0.0000)--(axis cs:7.0000,0.0000);\end{axis}\end{tikzpicture}\end{center}
\item Resol \fmm{\begin{cases}x+y\ge4\\ x+y\le1\end{cases}}

Les rectes \fmm{x+y=4} i \fmm{x+y=1} són para\lgem{}eles. Amb \fmm{O}: \fmm{0\ge4} és fals, així que el primer semiplà no conté l'origen; \fmm{0\le1\ \checkmark}, així que el segon sí.

Els dos semiplans no tenen cap punt en comú: el sistema no té solució.

\begin{center}\begin{tikzpicture}\begin{axis}[fmaxis,axis equal image,width=5.6cm,height=5.6cm,xmin=-3,xmax=7,ymin=-3,ymax=7,xtick={-2,2,4,6},ytick={-2,2,4,6},minor xtick={-3,-2,-1,0,1,2,3,4,5,6,7},minor ytick={-3,-2,-1,0,1,2,3,4,5,6,7},grid=both,minor grid style={fmmuted!15},major grid style={fmmuted!15},tick label style={font=\tiny}]\draw[fmA,line width=1.5pt] (axis cs:-4.0000,8.0000)--(axis cs:8.0000,-4.0000);\draw[fmB,line width=1.5pt] (axis cs:-4.0000,5.0000)--(axis cs:8.0000,-7.0000);\end{axis}\end{tikzpicture}\end{center}
\end{enumerate}
\end{fmex}

\Needspace{8\baselineskip}
\section{Resolució de problemes}\hypertarget{sec6}{}
\begin{fmprop}{Pautes per a resoldre problemes amb sistemes}
\begin{enumerate}
\item \textbf{Llegir} l'enunciat amb atenció i identificar les dades i el que es demana.
\item \textbf{Anomenar les incògnites} (\fmm{x,y,z}\dots{}) i indicar què representa cada una (i en quines unitats).
\item \textbf{Plantejar} el sistema: cada condició de l'enunciat és una equació o, si diu «com a màxim» o «com a mínim», una inequació.
\item \textbf{Resoldre'l.}
\item \textbf{Comprovar i interpretar} les solucions: descartar les que no tenen sentit en el context (valors negatius, no enters\dots{}) i respondre amb una frase.
\end{enumerate}
\end{fmprop}

\Needspace{5\baselineskip}
\begin{fmex}{Exemple resolt: problemes}
\begin{enumerate}[label=\textbf{\arabic*.},leftmargin=*]
\item En una botiga hi ha 40 articles entre samarretes de 12 € i pantalons de 25 €. Si el valor total és de 662 €, quants n'hi ha de cada tipus?

\fmm{x=} samarretes, \fmm{y=} pantalons: \fmm{\ \begin{cases}x+y=40\\ 12x+25y=662\end{cases}}

\fmm{12(40-y)+25y=662\Rightarrow13y=182\Rightarrow y=14,\ x=26}

Hi ha 26 samarretes i 14 pantalons.
\end{enumerate}
\end{fmex}

  
\end{document}
